% Generated draft: requires mathematical and source review.
\documentclass[11pt]{article}
\usepackage[margin=25mm]{geometry}
\usepackage{fontspec}
\setmainfont{Latin Modern Roman}
\usepackage{amsmath,amssymb,mathtools,unicode-math}
\setmathfont{Latin Modern Math}
\usepackage{enumitem,needspace,xurl,hyperref}
\setlength{\parindent}{0pt}
\setlength{\parskip}{5pt}
\setlength{\emergencystretch}{3em}
\newcommand{\N}{\mathbb{N}}
\newcommand{\Z}{\mathbb{Z}}
\newcommand{\Q}{\mathbb{Q}}
\newcommand{\R}{\mathbb{R}}
\newcommand{\C}{\mathbb{C}}
\newcommand{\F}{\mathbb{F}}
\newcommand{\A}{\mathbb{A}}
\newcommand{\PP}{\mathbb{P}}
\newcommand{\E}{\mathbb{E}}
\newcommand{\eps}{\varepsilon}
\newcommand{\dd}{\,\mathrm{d}}
\begin{document}
\section*{UnsolvedMath ID: 700807}
\textbf{Source problem code:} catalogue entry
\subsection{Definitions and mathematical statement}
Let \(X\) be a noncyclic free group, with no restriction in the supplied statement that its rank be finite. Let \(1\ne N\trianglelefteq X\) and let \(N\le T<X\). For subgroups \(A,B\le X\), define
\[
[A,B]=\langle a^{-1}b^{-1}ab\mid a\in A,\ b\in B\rangle.
\]
Here \(<\) denotes proper subgroup inclusion. In particular, \([T,N]\le [X,N]\).

The source asks whether, for every such triple \((X,N,T)\), the additional hypothesis that \(N\) is ``not maximal in \(X\)'' implies
\[
[T,N]<[X,N].
\]
Under the interpretation that maximality refers to all subgroups of \(X\), the question is precisely whether
\[
\left(\exists H\le X\;\;N<H<X\right)
\quad\Longrightarrow\quad
[T,N]\ne[X,N]
\]
holds for every noncyclic free group \(X\), every nontrivial normal subgroup \(N\) of \(X\), and every subgroup \(T\) satisfying \(N\le T<X\).

\emph{Unresolved qualification.} The supplied record does not specify whether ``maximal'' instead means maximal among proper normal subgroups. Under that interpretation, the additional hypothesis would be
\[
\exists H\trianglelefteq X\quad N<H<X.
\]
These two hypotheses need not coincide; the intended interpretation requires source review. The source permits \(T=N\) and does not require \(T\) to be normal in \(X\).
\subsection{Short English statement}
V. P. Shaptala asks whether a nontrivial normal subgroup \(N\) of a noncyclic free group \(X\), subject to the source's nonmaximality condition, satisfies \([T,N]<[X,N]\) for every proper subgroup \(T\) containing \(N\). The meaning of nonmaximality requires clarification. The supplied registry retains an open label and records an explicitly unverified AI claim of a negative answer; that claim is not treated here as an established result.
\subsection{Sources}
\begin{itemize}
\item Local UnsolvedMath registry, supplied complete source record: permanent ID \(700807\), exact source code \texttt{catalogue entry}, title \emph{Kourovka Notebook Problem 11.115}. Attribution: V. P. Shaptala, proposed in 1990.
\item Supplied primary-source attribution: \emph{The Kourovka Notebook (Unsolved Problems in Group Theory)}, 11th issue (1990), Problem 11.115; also identified in the supplied record with the 21st issue edition (2026), arXiv:1401.0300v48. The collection description names E. I. Khukhro and V. D. Mazurov as editors.
\item Supplied discussion and status-review note: the 21st issue (October 2026) marks the problem with \(\#\) and reports an explicitly unverified AI claim that the answer is negative, attributed to M. Kinyon and J. Urban, \emph{Preprint}, 2026: \url{https://kourovkanotebookorg.wordpress.com/wp-content/uploads/2026/09/kourovka-mathematics.pdf}. This is supplied attribution of an unverified claim, not verification of a counterexample or solution.
\item The supplied status review is dated \texttt{2026-10-16}, retains canonical status \texttt{open}, and classifies its evidence as \texttt{unrefereed\_claim}. This assessment is supplied dated curation, not independently verified literature assessment. No cited source was independently consulted for this draft.
\end{itemize}
\end{document}
